\section{Axis-wake inversion parity}\label{axis:sec:axis}
We use the complete-state and pre-update conventions of Section~2.

We study a nonempty turn word $w=w_0\cdots w_{P-1}$ whose infinite repetition
is executed from such an input. Its heading returns after $P$ updates, and its
\emph{full} displacement is a nonzero axis vector. A quarter-turn rotation
normalizes it to $d=(a,0)$ with $a>0$. The word need not be a least period.
The complete colouring need not translate: it retains all emitted wake and
unread debris. A finite transient before the repeated suffix is allowed,
since the field at its end is still finite.

Let $q_i$ be the pre-update position of phase $i$. Translation classes are
cells modulo $\Z(a,0)$, represented by $(x\bmod a,y)$. Let $O$ be the finite
set of classes containing an odd number of the $P$ phases, and let $g=|O|$.
Define the \emph{strict inversion count}
\begin{equation}\label{axis:eq:I}
 I(O)=\#\{(u,v)\in O^2:x_u<x_v,\ y_u>y_v\}.
\end{equation}
Equal rows or columns do not contribute. Equivalently, sort $O$ by increasing
$(y,x)$ and count strict inversions of the resulting $x$ list.

\begin{theorem}[Axis-wake obstruction]\label{axis:thm:main}
Every finite-support periodic realization with full displacement $(a,0)$,
$a>0$, has even $a$ and satisfies
\[
 \boxed{I(O)\equiv0\pmod2.}
\]
There is no bound on $P$, vertical span, $g$, even visit multiplicities,
initial black cardinality or initial coordinates.
\end{theorem}

The test is intrinsic to the cylindrical wake. Translating the coordinates,
changing the phase origin or changing representatives does not alter its
parity on the relevant domain.

\paragraph{Relation to earlier results.}
Jillhewar's stabilized-class criterion, signed mod-four residues and
special-axis lifted-rank factorization supply the underlying methodology
\cite{Jillhewar2026}. The non-axis wake character in Appendix~\ref{app:quadratic}
uses uniqueness of finite horizontal and vertical boundary antiderivatives.
That uniqueness fails on an axis cylinder: a whole horizontal row can be
added without changing any endpoint. The new step here is to prove that
these row cycles lie in the radical of the inversion polarization and to
eliminate their remaining scalar contribution using winding. A second step
reduces the resulting character to the endpoint order \eqref{axis:eq:I}.
For $a=2$ the result recovers the existing growth divisibility condition;
the new scope includes every larger even full axis displacement.

\subsection{Stable visits and the cylindrical parity graph}\label{axis:sec:stable}
For $q_i=b_v+k_i d$, the stable-time key of Section~\ref{sec:chronology}
is $\tau_i=i-Pk_i$. Lemma~\ref{lem:criterion} orders the class visits
$\Rr,\Ll,\Rr,\Ll,\ldots$ by these keys. Different phases have different keys,
since equality would imply equality modulo $P$.

The signed excess $D_v=\#\Rr-\#\Ll$ in a class is zero or one, and is exactly
its odd indicator. Hence the total excess is $g$. Heading restoration gives
$4\mid g$. It also makes $P$ even. Since horizontal and vertical moves
alternate and the net vertical displacement is zero, $P/2$ is even. Therefore
$4\mid P$, and the horizontal net displacement $a$ is even.

The product of checker parity with the sign distinguishing horizontal and
vertical incoming headings is invariant. Label vertical-arrival vertices
$Q$, horizontal-arrival vertices $C$, choosing the coordinate gauge in which
$Q$ has even $x+y$. Horizontal directed edges run from $Q$ to $C$ and vertical
edges from $C$ to $Q$. Each visit uses one edge of each axis. Consequently the
horizontal and vertical boundaries of the traversed-edge indicator modulo two
are both $D=\mathbf1_O$. Every row and every column of $O$ has even population.

Pair consecutive $y$ endpoints on each column to obtain the unique finite
vertical edge set with boundary $D$. On each row choose the horizontal seam
edge from column $a-1$ to column zero to be absent, and pair consecutive $x$
endpoints in $0,\ldots,a-1$. Call the union of these horizontal and vertical
interval edges $G_0(O)$. All edges retain their physical direction and port.
For $a=2$, the east and west edges between the same quotient vertices remain
distinct parallel edges.

\begin{lemma}[All possible parity lifts]\label{axis:lem:lifts}
The actual parity graph of any period with wake $O$ is
\begin{equation}\label{axis:eq:lifts}
 G=G_0(O)\xor\mathop{\triangle}_{y:h_y=1}R_y,
\end{equation}
where $R_y$ is the full horizontal row cycle and only finitely many $h_y$ are
nonzero. Moreover $\sum_yh_y\equiv1\pmod2$.
\end{lemma}
\begin{proof}
A finite boundary-free edge set on a vertical line is empty. On a horizontal
cycle it is either empty or the full cycle. Subtracting $G_0$ from the actual
edge indicator gives \eqref{axis:eq:lifts}. There is no assumption that the added
cycles are between the highest and lowest wake endpoints.

Lift the closed quotient walk to the plane. Its full displacement is $(a,0)$,
so its signed crossing number of the seam is one. Modulo two signs disappear.
The seam-zero graph contributes zero and each $R_y$ contributes one. This gives
the asserted parity. Even traversal multiplicities do not affect it.
\end{proof}

\subsection{The one-tour port sign}\label{axis:sec:sign}
Order the finite set of relevant vertices by any total order, and at each
vertex put the negative port before the positive port. Positive means movement
N or E, for both incoming and outgoing ports. For a finite directed parity
edge set $G$, let $E(G)$ be the inversion parity between the list of its edges
ordered by source port and the same edges ordered by target port. Let
$N_+(G)$ count its positive edges. Define
\begin{equation}\label{axis:eq:q}
 q(G;O)=E(G)+N_+(G)+|O\cap C|\pmod2.
\end{equation}
We suppress $O$ when it is fixed. The next factorization extends the
special-drift port argument in \cite{Jillhewar2026}; we include its derivation.

\begin{lemma}[Actual-period character]\label{axis:lem:actual}
The parity graph of every period under consideration satisfies $q(G)=1$.
\end{lemma}
\begin{proof}
At a vertex with $m$ stable-ranked events, record their positive incoming and
outgoing bits as $\alpha_k,\beta_k$, $0\le k<m$. Put
$A=\sum\alpha_k$, $B=\sum\beta_k$ and $p=A+B\pmod2$. For any binary list,
\begin{equation}\label{axis:eq:binary}
 \inv(\alpha)=A(m-1)+\sum_k k\alpha_k+\binom A2\pmod2.
\end{equation}
Indeed, a one at position $k$ sees $m-1-k$ later positions, from which its
later ones are subtracted; the latter count sums to $\binom A2$.

Write $K$ for the sum of the stable ranks of all left turns, modulo two, and
$I_{\rm port}$ for the sum of incoming and outgoing binary inversions.
At $Q$ a left turn means unequal port bits, whereas at $C$ it means equal
bits. Applying \eqref{axis:eq:binary} to both lists gives local corrections
$p(m-1)$ at $Q$ and $p(m-1)+\binom m2$ at $C$. The positive-edge binomial terms
cancel globally, because a directed edge has equal source and target counts.
At $C$ one also has
\[
 D=2(A+B)-4\sum_k\alpha_k\beta_k-m,\qquad m\equiv2p-D\pmod4.
\]
Using $D\in\{0,1\}$ and $m\equiv D\pmod2$ therefore gives
\begin{equation}\label{axis:eq:local}
 K=I_{\rm port}+L,\qquad
 L=\sum_Qp(D+1)+\sum_CD(p+1)\pmod2.
\end{equation}

A quotient edge of voltage $\nu$ changes stable time by $1-P\nu$. This same
constant applies to all its occurrences, including the period seam. Source
and target occurrences therefore match first to first. Sorting outgoing
ports, transferring edge blocks and unsorting incoming ports factors the
phase successor. That successor is one even-length $P$-cycle, even when the
chosen word is a power. Hence
\begin{equation}\label{axis:eq:tour}
 I_{\rm port}+E(G)=1\pmod2.
\end{equation}
Exchanging blocks of sizes $r,s$ contributes $rs$ modulo two, so the edge-block
sign depends only on edge-count parities, precisely as in the definition of
$E(G)$.

It remains to simplify $L$. In $G_0$, a paired horizontal or vertical interval
contributes positive-incidence parity $1+c(u)+c(v)$ at its endpoints, where
$c=0$ on $Q$ and $1$ on $C$. There are $g/2$ intervals of each axis. Each
endpoint checker bit is counted twice; thus $\sum_{v\in O}p_v=0$.
Adding $R_y$ changes $p_v$ once at every vertex on that row. The row population
of $O$ is even, so the same cancellation holds for every graph \eqref{axis:eq:lifts}.
Furthermore $\sum_Qp=N_+$: each positive edge has exactly one $Q$ endpoint.
Expanding \eqref{axis:eq:local} gives
\[
 L=N_+(G)+|O\cap C|\pmod2.
\]
Every left turn in a stable alternating class has odd rank, so
$K=\#\Ll=(P-g)/2\equiv0\pmod2$, since $4\mid P,g$.
Equations \eqref{axis:eq:local} and \eqref{axis:eq:tour} now give $q(G)=1$.
\end{proof}

\subsection{Eliminating unknown row circulations}\label{axis:sec:radical}
For a fixed edge universe and vertex order, $E$ is a quadratic polynomial over
$\F$ in the edge indicators: each inverted pair contributes the product of
its two indicators. Its polarization
\[
 \mathcal B(G,H)=E(G\xor H)+E(G)+E(H)
\]
is bilinear. The following lemma is the axis-specific cancellation.

\begin{lemma}[Row-cycle radical]\label{axis:lem:radical}
For every finite edge set $G$ and every full horizontal row cycle $R_y$,
\[
 \mathcal B(G,R_y)=0,\qquad E(R_y)+N_+(R_y)=1\pmod2.
\]
The first assertion does not require a boundary condition on $G$.
\end{lemma}
\begin{proof}
Every source and target degree of $R_y$ is even. Swapping vertices in the
common order therefore leaves $\mathcal B(G,R_y)$ unchanged. Place its row
vertices first, ordered by column, and use bilinearity to reduce to a single
edge $e$.

An edge off the row sees only complete source and target blocks of $R_y$,
each of size two, so its mixed inversion count is even. A vertical edge
touching the row has its row source at $C$, where $R_y$ has no source block,
or its row target at $Q$, where $R_y$ has no target block. The remaining
whole blocks again have size two. Finally a horizontal edge on the row is
itself in $R_y$. Its inversion count with the other row edges has parity equal
to the sum of its zero-based source and target indices. Both indices have
the parity of its port bit, so their sum is even. This also covers the two
parallel ports when $a=2$. Thus the polarization vanishes.

Put $a=2m$. With the row order just chosen, one sign of horizontal port
matches identically among $m$ slots and the other matches by a cyclic shift
among $m$ slots. Its sign is $m-1$ modulo two. There are $m$ positive edges.
Their sum is $2m-1\equiv1$.
\end{proof}

Since edge-number parity is linear, toggling a row cycle changes $q$ by one.
Lemma~\ref{axis:lem:lifts} and Lemma~\ref{axis:lem:actual} consequently imply
\begin{equation}\label{axis:eq:q0zero}
 q(G_0(O);O)=0.
\end{equation}
The entire unknown set of row circulations has disappeared. In particular,
no bound on its support or on the number of even visits was needed.

\subsection{Evaluating the character from endpoints}\label{axis:sec:endpoints}
The next statement is purely combinatorial, on a larger domain than that of
realizable wakes.

\begin{theorem}[Endpoint identity]\label{axis:thm:endpoint}
For even $a>0$ and any finite $O\subset\{0,\ldots,a-1\}\times\Z$ with even
populations on every row and column,
\begin{equation}\label{axis:eq:endpoint}
 q(G_0(O);O)=I(O)+|O|/2\pmod2.
\end{equation}
It holds for either checker/arrival-axis convention.
\end{theorem}
\begin{proof}
At every vertex the incoming and outgoing degree parities of $G_0$ are both
$D_v$. Exchanging vertices $v,w$ changes $E$ by
$D_vD_w+D_vD_w=0$. We may therefore put every $Q$ vertex before every $C$
vertex, using row-major order within each checker colour.

Every horizontal/vertical edge pair is inverted: its sources are in the order
$Q,C$ and its targets in the order $C,Q$. The parity of the horizontal edge
count is $|O\cap Q|$, and that of the vertical count is $|O\cap C|$.
The two endpoint populations have the same parity because $|O|$ is even.
Thus the mixed inversions cancel the checker term in \eqref{axis:eq:q}, leaving
\begin{equation}\label{axis:eq:axes}
 q(G_0)=(E_H+N_H)+(E_V+N_V).
\end{equation}

Consider a paired horizontal interval $[l,r]$. Its inversions are exactly its
internal $C$ vertices, where the two incoming edges reverse port order.
Its positive edges are counted by $Q$ vertices from $l$ inclusive to $r$
exclusive. Summing modulo two gives $c(r)$. Distinct horizontal intervals
have no mutual inversions in row-major order. Therefore
\[
 E_H+N_H=\sum_{\text{horizontal pairs}}c(\text{right endpoint}).
\]

To evaluate the vertical part, change to column-major order within each
checker colour. A paired vertical interval $[b,t]$ contributes
$1+c(t)$: its inversions are its internal $Q$ vertices and its positive
edges start at $C$ vertices from $b$ inclusive to $t$ exclusive. Distinct
paired vertical intervals have no mutual inversions in column-major order.
Changing the order back adds $J_{\rm same}$, the strict inversion count among
endpoint pairs of the same checker colour. This follows by block exchanges:
vertical source degrees at $C$ and target degrees at $Q$ have parity $D_v$.
Together with \eqref{axis:eq:axes}, this gives
\begin{equation}\label{axis:eq:ranks}
 q(G_0)=|O|/2+J_{\rm same}
 +\sum_{\text{horizontal pairs}}c(r)+\sum_{\text{vertical pairs}}c(t).
\end{equation}

Let $r(v)$ and $s(v)$ be the zero-based ranks of endpoints in row-major and
column-major order. Because every row and column has even population, paired
right endpoints are precisely the odd row-major ranks, and paired top
endpoints precisely the odd column-major ranks. The last two sums in
\eqref{axis:eq:ranks} equal $\sum_vc(v)(r(v)+s(v))$ modulo two.
The parity of an endpoint's rank change is the number of endpoints with which
it exchanges order. Each inverted pair $u,v$ therefore contributes
$c(u)+c(v)$ to this sum. The sum is exactly $J_{\rm mixed}$, the count of
inversions with unlike checker colours. Adding $J_{\rm same}$ proves
\eqref{axis:eq:endpoint}. The same argument works after exchanging the checker labels.
\end{proof}

\begin{proof}[Proof of Theorem~\ref{axis:thm:main}]
Section~\ref{axis:sec:stable} gives even $a$, even fiber populations and $4\mid g$.
Equation~\eqref{axis:eq:q0zero} and Theorem~\ref{axis:thm:endpoint} give
$0=I(O)+g/2=I(O)\pmod2$.
\end{proof}

\paragraph{Translation and phase invariance.}
Vertical translation does not change row order. To move the horizontal seam,
it suffices to move one column from the start of the column order to its end.
If this column has $n$ endpoints, the number of strict comparisons reversed is
\[
 n(g-n)-\sum_{v\text{ in that column}}(\#\text{endpoints in row }y_v-1).
\]
Both $n$ and $g$ are even; each row population is even, and there are an even
number $n$ of terms in the sum. The displayed number is even. Hence cyclic
column shifts preserve $I$ modulo two. Changing the phase origin translates
the odd classes: visits shifted across the period seam differ by full copies
of $d$ and remain in their classes. Normalizing representatives also leaves
the same set. These observations prove the asserted invariances without a
geometric cutoff.

\paragraph{Quadratic structure.}
On the vector space of finite indicators with even row and column populations,
$O\mapsto G_0(O)$ is linear: in each line or cut cycle, cumulative boundary
parity is a linear antiderivative. Thus $q(G_0(O);O)$ is a quadratic form.
Its polarization is bilinear. Addition here is symmetric difference of
indicators, not superposition of physical ant trajectories.

\subsection{A fast screen and excluded families}\label{axis:sec:families}
Given $g$ explicit class representatives in binary, reduce their horizontal
coordinates modulo $a$, deduplicate and sort by $(y,x)$. Merge-sort inversion
counting computes $I$ in $O(g\log(g+1))$ comparisons and arithmetic operations
and $O(g)$ space. Equal column values must be treated as noninversions.
The integers have polynomial bit size in the input coordinates and $\log(g+1)$.
The algorithm never expands a vertical gap.

For comparison with existing arithmetic screens we use the signed residue
conditions from \cite{Jillhewar2026}:
\begin{equation}\label{axis:eq:residues}
 \sum_{(x,y)\in O:\,x=r}(-1)^{x+y}\equiv2\pmod4,
 \qquad
 \sum_{(x,y)\in O:\,y=s}(-1)^{x+y}\equiv0\pmod4.
\end{equation}
They imply $g\ge2a$, and impose even populations on each fiber. For clarity,
their mod-four potential can be written explicitly. For arbitrary
$f:\Z\to\Z/4\Z$, choose $S$ with $S(x+1)-S(x)=f(x)+f(x+1)$ and put
\[
 F_N=-\chi S(x),\quad F_E=\chi(S(x)-f(x)),\quad
 F_S=-\chi(S(x)+2f(x)),\quad F_W=\chi(S(x)+f(x)),
\]
where $\chi=(-1)^{x+y}$. Direct substitution in the eight turn/heading cases
shows that $\sum_{z\in B}\chi(z)f(z_x)+F_h(p)$ is invariant modulo four.
For $a$-periodic $f$, its head monodromy along $(a,0)$ is $2\sum_{r=0}^{a-1}f(r)$.
The field increment is $\sum_{v\in O}\chi(v)f(v_x)$ by stable alternation.
Choosing residue indicators gives the first condition in \eqref{axis:eq:residues}.
Rotating the formula and using zero vertical displacement gives the second,
for arbitrary single-row indicators. This also specifies their exact-row scope.

\begin{theorem}[Staircase exclusions]\label{axis:thm:staircase}
For every even $a\ge2$ and every list of strictly increasing heights
$y_0<\cdots<y_{a-1}$ of a single parity, define
\[
 O_a=\bigcup_{j=0}^{a-2}\{(j,y_j),(j+1,y_j)\}
 \ \cup\ \{(0,y_{a-1}),(a-1,y_{a-1})\}.
\]
These $2a$ classes satisfy \eqref{axis:eq:residues}, $4\mid g$ and $g=2a$, but
no finite-support periodic trace of full drift $(a,0)$ has this wake.
\end{theorem}
\begin{proof}
Every column has two endpoints of the same checker sign, giving $\pm2\equiv2$
modulo four. Each row has two endpoints of opposite sign, giving zero.
The row-major column list is
\[
 0,1,\ 1,2,\ldots,\ a-2,a-1,\ 0,a-1.
\]
Only the penultimate zero causes strict inversions. Exactly $2a-3$ preceding
entries are positive. This is odd, contradicting Theorem~\ref{axis:thm:main}.
\end{proof}

For $a=4$, heights $0,10,20,30$ give eight endpoints and $I=5$.
For contrast, all four columns on rows zero and twenty give $I=6$ and pass
both the old conditions and the new screen. Passing establishes only
\texttt{NOT\_EXCLUDED}, not a realizable periodic word.

At $a=2$, even row population means that every occupied row contains both
classes. If there are $m$ occupied rows, $g=2m$ and $I=\binom m2$.
Since $4\mid g$, $m$ is even, so $I\equiv g/4\pmod2$. The theorem recovers
$8\mid g$, an already established two-column consequence of the attributed
rank argument. It does not claim a new proof of the separate condition
$P\equiv4\pmod8$.

\subsection{An exact rejection count in a factorial family}\label{axis:sec:count}
Fix even $a\ge4$ and $a$ strictly increasing same-parity heights. Let the
cycle edges on the columns be
\[
 E_j=\{j,j+1\}\ (0\le j<a-1),\qquad E_{a-1}=\{a-1,0\}.
\]
Assign these $a$ distinct edges in any order to the $a$ rows, placing endpoints
at the assigned columns. There are $a!$ distinct labelled wakes. Each column
still has degree two, and each row contains opposite checker signs, so all
these wakes satisfy the conditions of Theorem~\ref{axis:thm:staircase} before its
inversion test.

Let $\pi_j$ be the row rank assigned to edge $E_j$. Define
\[
 D_{\rm cyc}(\pi)=\sum_{j=0}^{a-1}\mathbf1_{\{\pi_j>\pi_{j+1\bmod a}\}}.
\]
Let $A(n,k)$ be the Eulerian number counting permutations of $n$ elements with
$k$ ordinary descents.

\begin{theorem}[Cycle-edge parity and counting]\label{axis:thm:count}
For this family,
\[
 I(O_\pi)\equiv D_{\rm cyc}(\pi)\pmod2.
\]
The number excluded by Theorem~\ref{axis:thm:main} is exactly
\begin{equation}\label{axis:eq:count}
 N_a=a\sum_{k\ {\rm even}} A(a-1,k).
\end{equation}
This counts labelled wake assignments rejected by the screen, not all
unrealizable wakes and not classes modulo symmetry.
\end{theorem}
\begin{proof}
Exchange two adjacent rows, carrying column pairs $E_i,E_j$.
Exactly $4-|E_i\cap E_j|$ unequal-column comparisons reverse. The inversion
parity therefore changes by $|E_i\cap E_j|$. For $a\ge4$ this is one precisely
when the two cycle edges are adjacent. Starting from the staircase order,
whose inversion count is $2a-3$, and sorting any row permutation by adjacent
exchanges gives
\[
 I(O_\pi)\equiv (2a-3)+
 \sum_{j=0}^{a-2}\mathbf1_{\{\pi_j>\pi_{j+1}\}}
 +\mathbf1_{\{\pi_0>\pi_{a-1}\}}.
\]
The final indicator is complementary to the closing cyclic descent. The odd
constant cancels, proving the parity formula.

Rotate the list $(\pi_0,\ldots,\pi_{a-1})$ until its smallest entry is first.
Every rotation orbit has $a$ members because its entries are distinct.
The first comparison is an ascent and the closing comparison a descent.
The remaining $a-1$ entries form an arbitrary permutation with, say, $k$
ordinary descents, so $D_{\rm cyc}=k+1$. There are $aA(a-1,k)$ such lists.
Odd cyclic descent count corresponds exactly to even $k$, giving \eqref{axis:eq:count}.
The recurrence
\[
 A(n,k)=(k+1)A(n-1,k)+(n-k)A(n-1,k-1),\qquad A(1,0)=1,
\]
follows by inserting the largest entry into a permutation: there are $k+1$
slots preserving $k$ descents, and $n-k$ slots increasing $k-1$ to $k$.
Thus the displayed count is effective for every $a$.
\end{proof}

\begin{center}
\begin{tabular}{rrrr}\toprule
Full displacement & Labelled wakes & Excluded & Not excluded\\\midrule
$(4,0)$ & 24 & 8 & 16\\
$(6,0)$ & 720 & 408 & 312\\
$(8,0)$ & 40320 & 19072 & 21248\\\bottomrule
\end{tabular}
\end{center}
The heights may have arbitrary gaps. Hence each rejected assignment supplies
an unbounded family of all-period exclusions.

